\section{Results}
\label{sec:results}

We report the four audit stages of \S\ref{sec:method}, followed by frozen
held-out evaluation of the final selected configurations. Aggregation,
per-head estimates and sensitivity analyses are detailed in
Appendices~\ref{app:stage1}--\ref{app:stage4}.

\subsection{RI poorly identifies causal contributors}
\label{sec:results-why}

\paragraph{What RI selects.}
Under our percentile rule, pooled RI selects two heads whose qualifying attention
is exclusively self- or previous-token attention; demonstrations contribute
over 99\% of their scores (Table~\ref{tab:pooled}), consistent with static
token preferences rather than retrieval. No eligible head passes the matched-target
test after Holm correction.
The test-only rule selects 59 heads, but several
are sensitive to the current token, target anchor or fact order
(Appendix~\ref{app:stage1}).

\begin{figure}[t]
\centering
\includegraphics[width=\columnwidth]{fig_ri_vs_importance}
\caption{RI versus exact signed Scope-P importance on 40 common pairs
(symmetric-log axis). Orange marks the 61 RI-selected heads; shading marks
$|I_h|<0.3$. (a) Pooled RI for all 1,024 heads. (b) First-token test-only
name gap for 167 supported heads.}
\label{fig:ri-vs-imp}
\end{figure}

\paragraph{Against the causal map.}
Pooled RI correlates weakly with signed importance ($\rho=0.13$), whereas
the test-only name gap correlates negatively ($\rho=-0.25$, or $-0.30$
within layer; Figure~\ref{fig:ri-vs-imp}). The two pooled-index heads have
near-zero effects, and 48 of the combined 61 selected heads have
$|I_h|<0.1$, less than 1\% of the clean--corrupted gap. Conversely, among
the 148 heads measured on all 178 pairs, only eight of the 25 heads with
$|I_h|\geq0.3$ are selected. Both selections are therefore incomplete,
and the pooled selection also meets our pre-registered misleading criterion
(Appendix~\ref{app:stage2}).

\paragraph{Why RI can miss retrieval.}
At the answer position the current token is always a colon, so when the visible
token IDs are unchanged, raw-embedding RI cannot adapt its name preferences
when mothers are reassigned (Appendix~\ref{app:stage1}). Copying a current
non-target name may also raise the name-control score and penalize a useful
head, as examined through the current-token and name-gap diagnostics in
Appendix~\ref{app:stage3}. These limitations concern the index, not the
model's capacity for semantic computation.

\subsection{The causal map}
\label{sec:results-map}

\paragraph{A few heads carry most of the effect.}
The gradient screen closely reproduces the exact ranking
($\rho=0.945$; Appendix~\ref{app:stage2}), and all effects reported below use
exact patching. Against a mean clean--corrupted gap of 11.2 logits, replacing
L17H1 changes the metric by 5.7 logits, about half the gap, while three further
heads each change it by roughly a quarter (Figure~\ref{fig:causal-map}).
Twenty-five heads reach $|I_h|\geq0.3$: 17 positive and eight negative, with
negative effects shifting the model toward the stated fact under corrupted-output
insertion. These are not
independent shares because head effects can interact. The same four heads lead
the exact common-subset scan, indicating that the leading heads are not an artefact of
the extension set (Appendix~\ref{app:stage2}).

\begin{figure}[t]
\centering
\includegraphics[width=\columnwidth]{fig_causal_map}
\caption{The causal map: signed Scope-P importance of the 25 strong heads
on 178 pairs, with family-bootstrap intervals; orange marks RI-selected heads.}
\label{fig:causal-map}
\end{figure}

\paragraph{Two positions, two kinds of head.}
Restricting the patch to the final colon splits the map by depth
(Figure~\ref{fig:scopeF-ratio}, Appendix~\ref{app:stage2}). In the Scope-F set,
measurable heads in layers 16--26 retain nearly all their
Scope-P effect under Scope F, indicating answer-position contributions, whereas
heads in layers 6--11 lose their effects and therefore write information at
earlier positions for later components to read; two layer-14--15 heads are
intermediate. At the colon, positive heads raise the donor-retrieved name and
negative heads lower it (Table~\ref{tab:scopeF-app});
\S\ref{sec:results-char} extends localization to all strong heads and tests
how their effects arise.

Together, these positional differences suggest a writer--reader organization
whose computations and connections are tested in
\S\S\ref{sec:results-char}--\ref{sec:results-routes}.

\subsection{What the important heads do}
\label{sec:results-char}
\paragraph{Writers and readers.}
Twenty of the 25 strong heads deliver nearly all their effect at the colon.
The strongest, L17H1, instead acts within the query fact, at mother tokens and
shared syntax positions (Figure~\ref{fig:profiles}). A controlled readout shows
that its write carries mother-identity information (Table~\ref{tab:readout-app}). L15H25 acts at the query
child's final token, but its written content remains unresolved. The strongest
writer thus shares the readers' layer band rather than belonging to the shallow
group of \S\ref{sec:results-map}.

\begin{figure}[t]
\centering
\includegraphics[width=\columnwidth]{fig_head_profiles}
\caption{Single-position effects of three fact-position heads by token role
(M: mother, C: child, oth.: other test facts, Q: question); mean and SD over
20 families.}
\label{fig:profiles}
\end{figure}

\paragraph{Question-guided selection, value-carried effects.}
At the colon, heads attend to mothers, children or the colon itself
(Figure~\ref{fig:colon-attention}). Mother- and child-attenders follow a changed
query to the newly relevant fact in all 20 families(Figure~\ref{fig:stage3-app}(a)); reordering or swapping
mothers has little effect on family-mean attention. Values carry most of the
leading heads' mean effects(Figure~\ref{fig:stage3-app}(b)). Together, these findings suggest question-guided
selection of positions whose content changes with the mother assignment.
L27H6's output directly favors the answer, whereas several strong heads have
output projections much smaller than their causal effects, motivating tests
of downstream computation. Even with contextual inputs instead of raw embeddings,
the RI name gap remains negatively correlated with causal importance
(Appendix~\ref{app:stage3}).

\paragraph{Copying does not determine causal role.}
Attended-name labels depend on the token anchor, and weight copying alone does not determine whether a head supports the correct answer. L18H18 has a positive causal effect ($I_h=+2.70$) despite a negative weight-copying score ($-0.08$); L30H18 has the highest weight-copying score ($+0.60$) but attends to a distractor and has a negative effect ($I_h=-1.06$; Table~\ref{tab:profiles-app}). Appendix~\ref{app:stage3} reports the event counts, anchor checks, control-name gaps, and synthetic-probe fingerprints.

\subsection{Measured routes and RI participation}
\label{sec:results-routes}

\paragraph{From profiles to connected paths.}
On 89 families, L17H1 reaches L18H18/L18H19 through V with effects of 3.40/2.07
logits; clamping its seven tested receivers removes 89.4\% of writer-union
restoration. Selectively releasing intermediate heads establishes the composed
writer--reader paths in Figure~\ref{fig:stage4-chains}, placing RI heads between
fact-site writers and later readers. L8H15 also has a direct V route despite
its weaker standalone effect; releasing MLP9 attenuates that positive route
(Appendix~\ref{app:s43-results}).

\begin{figure}[t]
\centering
\includegraphics[width=\columnwidth]{fig_stage4_chains}
\caption{Measured paths, 89 families. Orange: RI heads. Numbers: whole-path
noising/restoration effects. Intermediate heads are released jointly; Q/V denotes
the final receiver channel, with output evaluated at the colon.}
\label{fig:stage4-chains}
\end{figure}

\paragraph{Joint effects and weaker RI heads.}
Replacing L18H18 and L18H19 together reduces the answer margin by 5.63 logits.
Group-minus-member tests show that a head's contribution changes when its
partners are also replaced; singleton effects therefore do not predict group
behaviour. Separately, we test 31 RI-selected heads prioritized by earlier RI
rankings or causal effects (RI31). Excluding their eight strong members leaves
RI23. Replacing these 23 together changes family margins by 0.594 logits in
mean absolute magnitude, versus 0.215 for layer-matched non-RI groups. Thus
weaker RI heads matter collectively, although the signed effect reverses between
replacement baselines (Appendix~\ref{app:s42-residual-results}).

\paragraph{Additional participation depends on the background.}
The final experiment retains a 33-head set assembled from causal and route
evidence (C33), then adds 17 RI heads (C50). Neither preserves the task: on
20 families, C50 retains 11.5\% of the fact-swap gap, prefers the first chain's
mother in 85\% of cases and achieves 50\% candidate accuracy.

Within this partial background, six candidates with weak standalone effects,
chosen from existing RI rankings at different token anchors, are tested against
five previously measured routes, including those in Figure~\ref{fig:stage4-chains}
(rosters: Appendix~\ref{app:final-ri}). L9H16 and L1H27 improve the four-cell
binding contrast by 0.601 and 0.260 logits, positive in every family. Three
routes remain retained, but no candidate--route interaction passes the effect
rule, leaving the candidates' connections unresolved (Appendix~\ref{app:s44-results}).

On 87 held-out families, L9H16 and L1H27 reproduce the four-cell binding gains (0.585 and 0.261 logits), each positive in every family. Their original fact-swap effects are positive with mean replacement but nearly vanish with paired donors; RI31’s collective effect persists under both baselines, though the group includes eight already-strong heads. Under mean replacement, C50 retains only 12.8\% of the full fact-swap gap. Validation thus supports the binding effects in this partial background, not a sufficient circuit or new route assignments (Appendix~\ref{app:s44-validation}).
